\documentclass[12pt,a4paper]{article}
\usepackage[utf8]{inputenc}
\usepackage[spanish]{babel}
\usepackage{geometry}
\usepackage{listings}
\usepackage{xcolor}
\usepackage{parskip}
\usepackage{amsmath}
\usepackage{hyperref}

\geometry{margin=2.5cm}

\lstset{
  basicstyle=\ttfamily\small,
  keywordstyle=\color{blue}\bfseries,
  commentstyle=\color{gray},
  stringstyle=\color{orange},
  breaklines=true,
  frame=single,
  language=Java,
  numbers=left,
  numberstyle=\tiny\color{gray},
  literate={á}{{\'a}}1 {é}{{\'e}}1 {í}{{\'\i}}1 {ó}{{\'o}}1 {ú}{{\'u}}1 {ñ}{{\~n}}1 {Á}{{\'A}}1 {É}{{\'E}}1 {Ó}{{\'O}}1
}

\title{\textbf{Gato (Tic-Tac-Toe) con Negamax}\\
  \large Documentación Técnica}
\author{Inteligencia Artificial}
\date{}

\begin{document}
\maketitle
\tableofcontents
\newpage

\section{Descripción del Proyecto}
Implementación del juego \textbf{Gato} con inteligencia artificial basada en \textbf{Negamax}, una variante simplificada del algoritmo Minimax. A diferencia de la versión Minimax, este proyecto incluye \textbf{tres niveles de dificultad}: Fácil, Medio y Difícil.

\section{Lógica del Juego}

\subsection{Representación del estado}
El tablero de 3×3 se almacena en un arreglo unidimensional de 9 enteros. Cada celda puede ser \texttt{VACÍA}, \texttt{JUGADOR} o \texttt{CPU}.

\subsection{Niveles de dificultad}

\begin{itemize}
  \item \textbf{Fácil}: la CPU elige movimientos \textit{aleatorios} sin ninguna evaluación.
  \item \textbf{Medio}: la CPU aplica Negamax con \textbf{profundidad limitada} (2-3 niveles), lo que la hace imperfecta.
  \item \textbf{Difícil}: la CPU aplica Negamax \textbf{completo} sin límite de profundidad, siendo imbatible.
\end{itemize}

\section{Algoritmo Negamax}

\subsection{Fundamento matemático}
Negamax se basa en la simetría de los juegos de suma cero:
\[
  \text{minimax}(\text{MAX}, s) = -\text{minimax}(\text{MIN}, s)
\]
Esto permite escribir \textbf{una sola función} en lugar de dos (maximizador y minimizador).

\subsection{Pseudocódigo}
\begin{lstlisting}
función negamax(tablero, profundidad, jugadorActual):
    si hayGanador(tablero, jugadorActual):
        retornar +10 - profundidad
    si hayGanador(tablero, oponente(jugadorActual)):
        retornar -10 + profundidad
    si tableroLleno(tablero):
        retornar 0

    mejorPuntuacion = -infinito
    para cada celda vacía en tablero:
        tablero[celda] = jugadorActual
        puntuacion = -negamax(tablero, profundidad+1,
                              oponente(jugadorActual))
        tablero[celda] = vacío
        mejorPuntuacion = max(mejorPuntuacion, puntuacion)

    retornar mejorPuntuacion
\end{lstlisting}

El signo negativo antes de la llamada recursiva invierte la perspectiva: lo que es bueno para el oponente es malo para mí.

\subsection{Factor de profundidad}
La puntuación ajustada por profundidad (\texttt{+10 - profundidad}) hace que el algoritmo prefiera victorias rápidas y derrote lo más tarde posible.

\section{Arquitectura del Código}

El proyecto se divide en dos clases:

\begin{itemize}
  \item \textbf{\texttt{Clase\_Gato.java}}: contiene la lógica del juego, el algoritmo Negamax y las utilidades de evaluación del tablero.
  \item \textbf{\texttt{Gato\_Grafico.java}}: contiene la interfaz gráfica Swing, el menú de selección de dificultad y los listeners de eventos.
\end{itemize}

\section{Flujo de Ejecución}

\begin{enumerate}
  \item Al iniciar, se muestra un diálogo para seleccionar el nivel de dificultad.
  \item El jugador humano hace clic en una celda vacía.
  \item Se registra el movimiento y se verifica estado terminal.
  \item Dependiendo del nivel, la CPU:
    \begin{itemize}
      \item \textit{Fácil}: elige celda aleatoria.
      \item \textit{Medio}: ejecuta Negamax con profundidad máxima 3.
      \item \textit{Difícil}: ejecuta Negamax completo.
    \end{itemize}
  \item Se actualiza la GUI y se verifica nuevamente el estado terminal.
\end{enumerate}

\section{Comparativa Minimax vs Negamax}

\begin{center}
\begin{tabular}{|l|c|c|}
\hline
\textbf{Característica} & \textbf{Minimax} & \textbf{Negamax} \\
\hline
Número de funciones & 2 (max + min) & 1 (unificada) \\
Complejidad del código & Mayor & Menor \\
Resultado & Idéntico & Idéntico \\
Dificultades variables & Más complejo & Más sencillo \\
\hline
\end{tabular}
\end{center}

\end{document}
